Preprocessor helps designers to embed script like
code configuration directives into RTL code for reuse
oriented designs to improve code integrity. Conventionally,
designers are used to write one-time scripts (Perl/sed/awk/csh)
to preprocess their RTL for similar project usage. This methodology
is not clean enough. Verification engineers have to create additional
steps in \texttt{Makefile} to preprocess code.
\mhdl{}'s preprocessor uses \vlog{} style macro syntax, introduces
more flow control directives that help designers perform conditional and
repetitive configuration on RTL.

In addition to conventional \texttt{\`{}ifdef}, \texttt{\`{}ifndef},
\texttt{\`{}else}, \texttt{\`{}define} and \texttt{\`{}include}
macro directives, \mhdl{} introduces \texttt{\`{}for}, \texttt{\`{}if} and
\texttt{\`{}let} to enlarge the power of preprocessor (see following examples).

\autoref{lst:arb in pp} is a simple Round Robin Arbiter FSM implemented
in \mhdl{} with facilitating preprocessor. This arbiter can respond to
a configurable number of slaves, which is controlled by macro \texttt{SLV\_NUM}.
Once the \mhdl{} code is finished, various arbiters can be generated in \sv{} with
giving different values to \texttt{SLV\_NUM} when invoke \mhdlc{}.

\begin{lstlisting}[caption={Configurable Arbiter}, label={lst:arb in pp}]
fsm arb;

/*-\pp{for}-*/ (i=1; /*-\pp{i}-*/<=/*-\pp{SLV\_NUM}-*/; i++)
slave_grnt_/*-\pp{i}-*/ = 1'b0;
/*-\pp{endfor}-*/

/*-\pp{for}-*/ (i=1; /*-\pp{i}-*/<=/*-\pp{SLV\_NUM}-*/; i++)
 /*-\pp{let}-*/ j = /*-\pp{i}-*/ + 1

/*-\pp{if}-*/ /*-\pp{i}-*/ != /*-\pp{SLV\_NUM}-*/
SLAVE_/*-\pp{i}-*/: begin
   if ( slave_req_/*-\pp{i}-*/ ) begin
      slave_grnt_/*-\pp{i}-*/ = 1'b1;
      if ( slave_eof_/*-\pp{i}-*/ ) begin
         slave_grnt_/*-\pp{i}-*/ = 1'b0;
         goto SLAVE_/*-\pp{j}-*/;
      end
      else begin
         goto SLAVE_/*-\pp{i}-*/;
      end
   end
   else
     goto SLAVE_/*-\pp{j}-*/;
end

/*-\pp{else}-*/
SLAVE_/*-\pp{i}-*/: begin
   if ( slave_req_/*-\pp{i}-*/ ) begin
      slave_grnt_/*-\pp{i}-*/ = 1'b1;
      if ( slave_eof_/*-\pp{i}-*/ ) begin
         slave_grnt_/*-\pp{i}-*/ = 1'b0;
         goto SLAVE_1;
      end
      else
        goto SLAVE_/*-\pp{i}-*/;
   end
   else
     goto SLAVE_1;
end
/*-\pp{endif}-*/
/*-\pp{endfor}-*/
endfsm
\end{lstlisting}

In the listing above, the \texttt{\`{}for} directive repetitively writes code with
slight differences. Default values of FSM output are set in the first \texttt{\`{}for} block.
The second \texttt{\`{}for} block writes slave handling code, one state for each slave.
Since it is a Round Robin arbiter, every state performs the same task: grant slave access if
there is a request, move to next slave when current one has no request or transaction is done,
roll back to the first slave when an arbitration round finishes.
The \texttt{\`{}if}/\texttt{\`{}else}/\texttt{\`{}endif} block checks whether the current state
is for the last slave.
The \texttt{\`{}let} directive performs arithmetic operations and calculates the value of
\texttt{j}, which is the number of the next slave.

\autoref{lst:arb in sv} is the generated \sv{} with \texttt{SLV\_NUM} set to $4$.

\begin{lstlisting}[caption={Generated 4-slave arbiter \sv{} code}, label={lst:arb in sv}]
module arbiter (
  clock, reset_n,
  slave_eof_1, slave_eof_2, slave_eof_3, slave_eof_4,
  slave_grnt_1, slave_grnt_2, slave_grnt_3, slave_grnt_4,
  slave_req_1, slave_req_2, slave_req_3, slave_req_4);

// ... declarations omitted for brevity ...

always_comb begin
  slave_grnt_1 = 1'b0; slave_grnt_2 = 1'b0;
  slave_grnt_3 = 1'b0; slave_grnt_4 = 1'b0;
  unique case ( 1'b1 )
    arb_cs[_SLAVE_1_] : begin
      if ( slave_req_1 ) begin
        slave_grnt_1 = 1'b1;
        if ( slave_eof_1 ) begin
          slave_grnt_1 = 1'b0; arb_ns = SLAVE_2;
        end else begin arb_ns = SLAVE_1; end
      end else begin arb_ns = SLAVE_2; end
    end
    // ... SLAVE_2, SLAVE_3 states follow the same pattern ...
    arb_cs[_SLAVE_4_] : begin
      if ( slave_req_4 ) begin
        slave_grnt_4 = 1'b1;
        if ( slave_eof_4 ) begin
          slave_grnt_4 = 1'b0; arb_ns = SLAVE_1;  // roll back
        end else begin arb_ns = SLAVE_4; end
      end else begin arb_ns = SLAVE_1; end
    end
    default: arb_ns = 4'hX;
  endcase
end
endmodule
\end{lstlisting}


\texttt{\`{}let} directive supports sufficient operators and a bunch of commonly used functions:
\begin{itemize}
\item Numeric operators: addition (+), subtraction ($-$),  multiplication ($*$), division ($/$),
modulus (\%), power ($**$).
\item Logical operators: logical AND (\&\&), logical OR ($||$), logical NOT ($!$).
\item Bit operators: bitwise XOR (\^{}), bitwise AND  (\&),  bitwise OR ($|$),
shift right ($>>$), shift left ($<<$).
\item Functions: log 2 (LOG2()), round up (CEIL()), round down (FLOOR()),
round to nearest value (ROUND()),  max of two numbers (MAX()),
min of two numbers (MIN()), odd (ODD()), even (EVEN()), absolute value (ABS()).
\end{itemize}

\begin{lstlisting}[caption={\texttt{\`{}let} usage examples}, label={lst:let example}]
/*-\pp{define}-*/ x 2

// NOTE: /*-\pp{let}-*/ need "="
/*-\pp{let}-*/ y = /*-\pp{x}-*/ ** 10   // y is 1024

/*-\pp{let}-*/ z = LOG2(/*-\pp{y}-*/)  // z is 10
/*-\pp{let}-*/ a = LOG2(/*-\pp{z}-*/)  // a is 3.321928
/*-\pp{let}-*/ c = CEIL(/*-\pp{a}-*/)  // c is 4
/*-\pp{let}-*/ f = FLOOR(/*-\pp{a}-*/) // f is 3
/*-\pp{let}-*/ r = ROUND(/*-\pp{a}-*/) // r is 3

// concatenate macro value with other strings using "::"
assign cat_/*-\pp{f}-*/::k = 1'b0;  // expand to ``assign cat_3k = 1'b0;''
assign cat_/*-\pp{fk}-*/ = 1'b0;    // /*-\textcolor{red}{Error!}-*/ macro ``fk'' undefined
\end{lstlisting}
\autoref{lst:let example} demonstrates some usages of \texttt{\`{}let} directive.
Operators and functions supported by \texttt{\`{}let} directive constructs macro variables
to \emph{macro expression}, such as
\texttt{\`{}a + \`{}b},
\texttt{ODD(\`{}f)},
\texttt{\`{}a >> 2},
\texttt{\`{}f <= \`{}d \&\& \`{}f != 0},
which can further be evaluated by \texttt{\`{}if} directive for flow controlling,
just as shown in \autoref{lst:arb in pp}.
